% =========================================================================
% Publisher Reliability Tool -- ECIR 2027 demo paper
% Springer LNCS, single-blind review (NOT anonymous).
% Format limit: max 4 pages of content, extra pages for references only.
% Submission deadline: 2 November 2026.
%
% Build:  pdflatex main && bibtex main && pdflatex main && pdflatex main
%
% Search this file for "TODO:" -- every open item is marked that way.
% =========================================================================
\documentclass[runningheads]{llncs}

\usepackage[T1]{fontenc}
\usepackage{newtxtext}
\usepackage[varvw]{newtxmath}
\usepackage{graphicx}
\usepackage{booktabs}
\usepackage{xcolor}
\usepackage{hyperref}
\usepackage{xurl}   % lets long URLs break across lines
\renewcommand\UrlFont{\color{blue}\rmfamily}
\urlstyle{rm}

% ---- Placeholder figure box -------------------------------------------
% Lets the paper compile before the screenshot exists. Replace the
% \figplaceholder{...} call with the \includegraphics line shown beside it.
\newcommand{\figplaceholder}[2]{%
  \fbox{\begin{minipage}[c][#1][c]{\dimexpr\linewidth-2\fboxsep-2\fboxrule\relax}
    \centering\small\ttfamily TODO: screenshot -- #2
  \end{minipage}}%
}

\begin{document}

\title{Publisher Reliability Tool: Interactive, Leakage-Safe\\
       Exploration of Article-to-Publisher Reliability Predictions}
\titlerunning{Publisher Reliability Tool}

\author{%
John Bianchi\inst{1}\orcidID{0009-0006-2582-1480} \and
Manuel Pratelli\inst{2}\orcidID{0000-0002-9978-791X} \and
Fabio Pinelli\inst{1}\orcidID{0000-0003-1058-6917} \and
Marinella Petrocchi\inst{2,1}\orcidID{0000-0003-0591-877X}}
\authorrunning{J. Bianchi et al.}

\institute{%
IMT School for Advanced Studies Lucca, 55100 Lucca, Italy\\
\email{\{john.bianchi,fabio.pinelli\}@imtlucca.it}
\and
Institute of Informatics and Telematics, National Research Council (IIT-CNR),
56124 Pisa, Italy\\
\email{\{manuel.pratelli,marinella.petrocchi\}@iit.cnr.it}}

\maketitle

% =========================================================================
\begin{abstract}
Automatic assessment of news source reliability is usually reported as
aggregate offline numbers, which makes it hard to inspect why a particular
publisher received a particular verdict. We present the \emph{Publisher
Reliability Tool} (PRT), an open-source local web application that turns a
five-fold, publisher-disjoint article classification study into an interactive,
auditable workspace. PRT ships a prediction-only release of 38,854
article-level runs over 19,411 political news articles and 838 publishers,
produced by BERT and RoBERTa classifiers trained on NewsGuard reliability
bands. Users browse stored predictions with complete per-class probabilities,
classify new article URLs on CPU, and aggregate 2--50 article runs into a
publisher-level verdict under three documented formulas. Two properties
distinguish PRT from a notebook: a server-side cross-validation leakage guard
that refuses to evaluate an article with a checkpoint trained on it, and
content-addressed model identities that keep every result reproducible from
plain, inspectable CSV. The licensed provider ratings themselves are never
redistributed.
\keywords{News source reliability \and Interactive demonstration \and
Prediction provenance \and Cross-validation leakage \and Reproducibility}
\end{abstract}

% =========================================================================
\section{Introduction}
\label{sec:intro}

Estimating how reliable a news \emph{publisher} is has become a recurring task
in information retrieval~\cite{nakov2024survey,srba2025survey}. A common design
classifies individual articles and aggregates those decisions to the
publisher~\cite{burdisso2024reliability}, and reports the outcome as a table of
offline scores: the reader learns how a model performed on average, but cannot
ask what it said about a specific publisher, which articles drove that verdict,
or what would change under a different aggregation rule. Since a single number
attached to a news outlet invites over-interpretation when the evidence behind
it is invisible, an interactive tool is what makes ``\emph{this is a model
prediction, conditioned on these exact inputs}'' inspectable rather than
asserted -- and it introduces a failure mode offline evaluation avoids by
construction: once a checkpoint can be aimed at an arbitrary article, nothing
prevents fold~$N$ meeting an article it was trained
on~\cite{verhoeven2025yesterday}.

We present the \textbf{Publisher Reliability Tool} (PRT), the reference
implementation accompanying our study on aggregating language model predictions
for news source reliability~\cite{bianchi2026articles}. It contributes an
offline-capable workspace over a prediction-only release of 38,854 article runs
with full probability vectors; a \emph{leakage-safe evaluation model} whose
offered checkpoints come from the intersection of installed artifacts, stored
coverage and held-out fold membership; and a content-addressed model identity
over seven human-readable CSV ledgers.

% =========================================================================
\section{Background and Related Work}
\label{sec:related}

Professional rating systems such as NewsGuard, the Global Disinformation Index
and Media Bias/Fact Check assess outlets at the \emph{publisher} level and agree
strongly, which makes them a common source of
supervision~\cite{lin2023correspondence,pratelli2022structured}. Computational
work estimates source reliability from text, interaction signals, or language
models used directly as
raters~\cite{burdisso2024reliability,yang2025accuracy,srba2025survey}; what
varies, and is rarely exposed, is the aggregation rule and the article sample
it consumes, which PRT turns into selectable parameters. Generalisation across
publishers is the fragile part~\cite{azizov2024safari,verhoeven2025yesterday},
which PRT keeps visible by treating fold membership as a service-enforced
safety property. Interactive publisher-reliability tools have appeared at ECIR
before~\cite{pratelli2025tropic}; PRT differs in publishing \emph{predictions
only}, so the licensed ratings are never redistributed.

% =========================================================================
\section{System Overview}
\label{sec:system}

\subsubsection{Task and data.}
PRT models a two-level task: a classifier maps one English article body to an
ordered class in $\{0,\dots,4\}$ with a five-component softmax vector, and an
aggregation function combines $2$--$50$ runs from \emph{one exact model} and
\emph{one normalised publisher} into a single class. The classes are the five
NewsGuard reliability bands in ascending order: \texttt{Class~0} is a score of
0--39 (proceed with maximum caution), then 40--59, 60--74, 75--99, and
\texttt{Class~4} is 100. Each article inherits its publisher's band, so
supervision is weak by construction~\cite{bianchi2026articles}. The corpus was
scraped with Newspaper3k from the homepages of NewsGuard-rated outlets tagged
\emph{Political News and Commentary}: English, paywall-free, capped per outlet,
at least 100 words.

% TODO: the source study reports 439 scraped domains (Table I) while the shipped
% TODO: release contains 838 distinct publisher hostnames over 19,429 rows, and
% TODO: its largest outlet has 294 articles against a stated cap of 200.
% TODO: Reconcile the two before submission and adjust the numbers below.
Crucially, PRT redistributes \emph{model outputs only}: the licensed scores,
bands and metadata are stripped at import and never persisted, logged or
rendered, and a class is always framed as a prediction. The
release contains 19,429 rows which, after canonical URL normalisation, yield
19,411 articles, 38,854 immutable runs and $10$ historical model/fold
identities over 838 publishers, $597$ with at least two articles (median
$15$).

\subsubsection{Architecture.}
PRT is a single Python~3.12 process that owns one data directory, serves the UI
and REST API from FastAPI on \texttt{127.0.0.1}, and runs one FIFO worker.
Frontend and API share the same request types and service functions, so the
browser cannot reach behaviour the API does not expose. All state lives in
seven CSV ledgers written by a single \texttt{flock}-guarded writer, three of
them append-only and immutable; malformed data fails closed before any HTTP
port is served. Articles and publishers are not stored -- they are views
derived at query time from prediction runs. Every local inference is mirrored
into the release CSV as a \texttt{user\_evaluation} row while original rows keep
\texttt{dataset\_original}; the release digest covers original rows only, so
user evaluations never alter release identity.

\subsubsection{Models and reproducible identity.}
The demonstrated path is CPU-only. BERT~\cite{devlin2019bert}
(\texttt{bert-base-uncased}) and RoBERTa~\cite{liu2020roberta}
(\texttt{roberta-large}) are fully fine-tuned five-label classifiers truncated
to 256 tokens, loaded with strict key and shape checking and gated by a frozen
CPU reference fixture. The loader interface is deliberately extensible: the
study's Llama~3~8B~\cite{llama2024herd} and
Mistral~24B~\cite{jiang2024mixtral} QLoRA checkpoints are authenticated against
a packaged manifest of exact filenames, byte sizes and SHA-256 digests, and
users may register their own five-class Transformers as a constrained ZIP
bundle -- inspectable, but outside the live demonstration since each fold needs
CUDA and several gigabytes. No path executes artifact-supplied code:
\texttt{auto\_map}, \texttt{trust\_remote\_code}, pickle weights and unsafe
archive paths are rejected.

A runnable model identity is the SHA-256 of canonical JSON over every
output-relevant setting -- artifact digest, loader recipe, base and tokenizer
revisions, class order, token limit, padding policy, dtype, quantisation -- but
not filesystem location; it is re-verified before loading, and imported
predictions get a separate, non-runnable \emph{historical} identity.

\subsubsection{Leakage-safe evaluation.}
The study enforces a publisher-disjoint protocol: fold $N$ means the stored
prediction was produced for held-out test fold $N$, so that checkpoint trained
on the other four. PRT makes this a service-layer invariant: checkpoint fold
$N$ may evaluate a known article in test fold $N$; a known article in any other
fold is rejected with the stable error
\texttt{TRAINING\_DATA\_LEAKAGE}, and publisher aggregation excludes every known
article outside its held-out fold. The $16$ articles normalising into more than one fold are excluded
entirely, and the guard covers single articles, explicit lists and publisher
candidates alike. PRT is equally careful about what this does \emph{not} prove:
absence from the registry is no evidence that an external article was outside
every training corpus.

\subsubsection{Retrieval and aggregation.}
Classifying a new URL uses an SSRF-hardened client that validates every
redirect and rejects non-public addresses; the HTML is parsed once by a single
English-configured extractor, and the text must reach 200 characters and 30
tokens and be detected as English before tokenisation, while raw HTML and
authors are discarded. Three versioned aggregation methods each require
$2$--$50$ runs from one model: \texttt{majority\_vote} -- the rule used in the
study -- \texttt{ordinal\_mean} (class $=\lfloor\mu+0.5\rfloor$) and
\texttt{mean\_probabilities}. Every evaluation records the ordered article and
run identifiers it consumed, so a later run can never change an earlier one.

% =========================================================================
\section{Demonstration}
\label{sec:demo}

The demonstration runs on an AMD Ryzen~7~2700U laptop (4 cores, 16\,GB RAM,
Ubuntu 26.04, no GPU) and needs no network for browsing or aggregation.
\emph{Browsing} an article lists every run with its probability vector, family
and fold, run identifier and origin. The visitor then pastes a public English
article URL: before submission PRT reports whether it is known, whether
inference is required, whether a matching checkpoint exists, and whether it is
blocked for leakage (Fig.~\ref{fig:evaluate}), so the selector is derived
rather than fixed. \emph{Reuse} returns the stored run without touching
the network; \emph{recompute} classifies the page into a new immutable run in
5--7\,s end to end (warm model time 0.9\,s for BERT, 2.9\,s for
RoBERTa-large). In \emph{publisher} mode the visitor picks a checkpoint, a
method and a count; PRT aggregates only stored, leakage-safe runs for that
hostname, never crawls, and re-running under another method makes the verdict's
sensitivity to the decision rule visible. We close by opening the CSV ledgers
beside the browser to show the new run and its mirrored release row.

\begin{figure}[t]
\centering
\figplaceholder{4.0cm}{Evaluate page -- URL field, derived model/fold selector with availability explanations, result card with probability bars, model/fold, run ID, origin}
% TODO: replace the \figplaceholder line above with the real screenshot.
% TODO: The 4.0cm box is sized for a TIGHTLY CROPPED screenshot. A taller
% TODO: image will push the paper past the 4-page content limit, so crop the
% TODO: Evaluate panel rather than capturing the whole browser window:
% \includegraphics[width=\textwidth]{figures/evaluate.png}
\Description{Screenshot of the Evaluate page: a URL input field, a derived
model and fold selector listing only leakage-safe checkpoints with a short
explanation for each unavailable option, and a result card showing the
predicted class, five probability bars, the exact model and fold, the
prediction run identifier and whether the run was stored or newly computed.}
\caption{Deriving the available checkpoints on the \emph{Evaluate} page.}
\label{fig:evaluate}
\end{figure}

% =========================================================================
\section{Availability}
\label{sec:availability}

% TODO: record the screencast and replace the placeholder URL below. ECIR
% TODO: requires a working online demo URL *or* a video of the main features;
% TODO: PRT is loopback-only by design, so the video is the route to take.
% TODO: add an archived release DOI (e.g. Zenodo) if one is minted.
A screencast of the walkthrough in Section~\ref{sec:demo} is at
\url{https://TODO-video-url}; the source code is under Apache-2.0 at
\url{https://github.com/bianchi-john/publisher-reliability-tool} with the
dataset release and its scientific and storage contracts. PRT starts with one
native command or one Compose service on \texttt{127.0.0.1:8000};
checkpoints and supplementary material~\cite{bianchi2026supplementary} are on
OSF\footnote{\url{https://osf.io/r9atz/overview?view_only=e4bda170a3e74ca3ae245475d4486d74}}.
Project-owned outputs use a limited
CC0 dedication; third-party URLs, names, ratings and weights are excluded.

% =========================================================================
\section{Conclusion and Future Work}
\label{sec:conclusion}

PRT turns a publisher-disjoint reliability study into a workspace where every
verdict traces to the checkpoint, articles and rule behind it; its leakage
guard and content-addressed identity generalise to any tool exposing
cross-validated classifiers. The release holds encoder predictions only and
aggregation is limited to three unweighted rules; a local evaluation mode would
let users with their own ratings measure the accuracy PRT does not report.

% =========================================================================
\begin{credits}
\subsubsection{\ackname}
Partially supported by the IIT-CNR internal project SIMULARE.
% TODO: add any further grant numbers, or adjust if the demo paper is funded
% TODO: differently from the underlying study.

\subsubsection{\discintname}
The authors have no competing interests to declare.
\end{credits}

\bibliographystyle{splncs04}
\bibliography{references}

\end{document}
